\section{Literature Review}
\begin{frame}{Beyond Precision and Recall: Measuring Search Engine Consistency Using Rank Stability}
    \scriptsize
    \vspace{-0.8em}
    
    \begin{table}[h!]
    \centering
    \renewcommand{\arraystretch}{1.5} % increases row height slightly more
    \setlength{\tabcolsep}{2.5pt}     % reduces horizontal padding a bit more
    
    \begin{tabular}{|p{2.6cm}|p{4.3cm}|p{4.3cm}|}
    \hline
    \textbf{Citation} & \textbf{Summary} & \textbf{Advantages \& Disadvantages} \\
    \hline
    \tiny\textit{[4] K. Batri \textit{et al.}, IEEE Access, vol. 13, pp. 92242–92259, 2025.  
    DOI: 10.1109/ACCESS.2025.3571184.}
    &
    Introduces \textbf{Rmeasure}, a framework to evaluate search engine consistency by measuring rank stability across repeated queries. It defines \textbf{Roverlap} and \textbf{Rnon-overlap} metrics and uses psychophysical modeling to account for user perception of rank changes. Experimental results show Bing demonstrates superior ranking stability compared to Google.
    &
    \textbf{Advantages:}
    \begin{itemize}\itemsep0.25em
    \item Evaluates overlapping/non-overlapping results.
    \item Strong statistical validation.
    \item Relevant to real-world query variability.
    \end{itemize}
    
    \textbf{Disadvantages:}
    \begin{itemize}\itemsep0.25em
    \item Small dataset (10 queries, 2 engines).
    \item Arbitrary parameters.
    \item Ignores personalization and time-based effects.
    \end{itemize} \\
    \hline
    \end{tabular}
    \end{table} 
\end{frame}